\section{Related Work}
\label{sec:related_work}

\subsection{Automatic Prompt and Program Optimization}

The strong sensitivity of LLMs to natural-language instructions has motivated a broad line of work on automated prompt optimization.
Early approaches such as Automatic Prompt Engineer generate and select instruction candidates using task-level evaluation
\citep{zhou2022ape,pryzant2023automatic}.
OPRO casts optimization itself as a language modeling problem, where an LLM proposes new solutions conditioned on previously evaluated candidates
\citep{yang2024opro}.
PromptBreeder adopts an evolutionary perspective and jointly evolves task prompts and the mutation prompts used to modify them
\citep{fernando2024promptbreeder}.
Other approaches optimize more structured language-model programs.
DSPy treats prompts and demonstrations as parameters of declarative LM pipelines and compiles them against task metrics
\citep{khattab2024dspy},
while TextGrad propagates textual feedback through compound AI systems in a manner analogous to automatic differentiation
\citep{yuksekgonul2024textgrad,agrawal2025textgrad}.
More recently, SkillOpt treats external agent skills as optimizable textual parameters while keeping the target model frozen
\citep{yang2026skillopt}.

These approaches establish that natural-language instructions can be optimized systematically without relying exclusively on manual prompt engineering. They also expose a recurring limitation: most optimization signals are collected at the level of complete trials, so the optimizer receives little information about when a prompt becomes harmful within a trajectory.
However, most of them optimize prompts or textual parameters using an offline development process and subsequently deploy the selected configuration.
MASOpt instead focuses on prompt adaptation \emph{during} the execution of a new multi-agent task.


\subsection{System-Level Prompt Optimization for Multi-Agent Systems}

Prompt optimization becomes more difficult in MAS because prompts interact through the workflow topology.
Optimizing one agent in isolation may improve its local behavior while changing the information distribution received by downstream agents.
Recent work has therefore moved from independent prompt optimization toward system-level coordination.
MAPRO formulates multi-agent prompt optimization as maximum a posteriori inference and uses topology-aware downstream feedback to selectively refine agent prompts
\citep{mapro2026}.
MASPOB formulates fixed-topology MAS prompt optimization as a budget-constrained black-box optimization problem.
It combines bandit exploration with a graph neural network surrogate and coordinate ascent to model topology-induced prompt interactions under a limited evaluation budget
\citep{hong2026maspob}.
MASPO instead evaluates prompts according to both local behavior and their contribution to downstream system success, and searches for coordinated prompt configurations across agents
\citep{wang2026maspo}.

These methods demonstrate the importance of optimizing MAS prompts as a coupled system rather than as independent components.
Their optimization target, however, remains a prompt configuration selected before deployment.
MASOpt addresses a complementary problem in which the prompt state itself remains adaptable after execution has begun.


\subsection{Runtime Feedback and Online Adaptation for Language Agents}

A related line of research studies how language models and agents can improve from feedback available at test or runtime.
Self-Refine alternates between feedback generation and output refinement, showing that an LLM can iteratively improve its own generations without additional parameter training
\citep{madaan2023selfrefine,pan2024critic}.
Reflexion converts environmental or task feedback into verbal reflections and stores them in episodic memory to improve subsequent agent trials
\citep{shinn2023reflexion,yao2023react}.
Recent memory-based approaches further explore persistent runtime adaptation.
For example, MemRL learns utilities over episodic memories from environmental feedback and uses them to improve future agent behavior while keeping the base model fixed
\citep{zhang2026memrl}.
These studies highlight the value of feedback, reflection, and non-parametric adaptation for language agents.

MASOpt differs in both its optimization target and temporal scope.
Rather than refining a single generated output or updating a persistent agent memory across tasks, MASOpt directly modifies the prompts of interacting agents inside the current MAS execution.
The effect of each modification is explicitly verified through the subsequent system response, making online prompt adaptation itself a closed-loop optimization process.
